\documentclass[10pt,a4paper]{amsart}
\usepackage{amsmath,amssymb,amsthm,graphicx}
\usepackage{hyperref}
\numberwithin{equation}{section}

\newtheorem{Theorem}{Theorem}[section]
\newtheorem{thm}{Theorem}[section]
\newtheorem*{thm*}{Theorem}
\newtheorem{cor}[thm]{Corollary}
\newtheorem{lem}[thm]{Lemma}
\newtheorem{prop}[thm]{Proposition}

\theoremstyle{definition}
\newtheorem{defin}[thm]{Definition}
\newtheorem{Note}[thm]{Note}
\newtheorem{Example}[thm]{Example}
\newtheorem{rem}[thm]{Remark}
\newtheorem{Question}[thm]{Question}

\def\C{{\bf C}}
\def\N{{\bf N}}
\def\Q{{\bf Q}}
\def\R{{\bf R}}
\def\Z{{\bf Z}}
\def\E{{\mathbb E}}
\def\V{{\bf V}}
\def\J{{\mathbb J}}
\def\W{{\bf W}}
%
\def\G{{\mathcal G}}
\def\T{{\mathcal T}}
\def\I{{\mathcal I}}
\def\D{{\mathcal D}}
\def\A{{\mathcal A}}
%
\newcommand{\cO}{{\mathcal O}}
\newcommand{\cC}{{\mathcal C}}
\newcommand{\cE}{{\mathcal E}}
\newcommand{\cL}{{\mathcal L}}
\newcommand{\cV}{{\mathcal V}}
\newcommand{\cW}{{\mathcal W}}
\newcommand{\cS}{{\mathcal S}}
\newcommand{\BR}{{\bf R}}

%
\newcommand{\lieg}{\mathfrak{g}}
\newcommand{\lieh}{\mathfrak{h}}
\def\g{{\mathfrak g}}
\def\h{{\mathfrak h}}
\def\k{{\mathfrak k}}
\def\p{{\mathfrak p}}
\def\q{{\mathfrak q}}
\def\a{{\mathfrak a}}
\def\b{{\mathfrak b}}
\def\m{{\mathfrak m}}
\def\n{{\mathfrak n}}
\def\z{{\mathfrak z}}
\def\t{{\mathfrak t}}
\def\lieo{{\mathfrak o}}
\newcommand{\so}{\mathfrak{so}}
\newcommand{\gl}{{\mathfrak{gl}}}
\newcommand{\sgl}{{\mathfrak{sl}}}
\newcommand{\hol}{\mathfrak{h}\mathfrak{o}\mathfrak{l}}
\newcommand{\Aut}{\operatorname{Aut}}
\newcommand{\Ad}{\operatorname{Ad}}
\newcommand{\ad}{\operatorname{ad}}
\newcommand{\GL}{{\operatorname{GL}}}
\newcommand{\SU}{{\rm SU}}
\newcommand{\SO}{{\rm SO}}

\title{Gromov-Zimmer embedding for parallel tractor sections with invariant connections}
\author{Karin Melnick, Katharina Neusser}
\date{doi:10.3842/SIGMA.2026.061; CC BY 4.0}
\begin{document}
\maketitle
\noindent\textit{Allowed standard prerequisites (editorial).}\begin{itemize}
\item Standard Cartan/tractor geometry and algebraic density. Original bibliography and complete source-local definitions retained.
\end{itemize}\medskip
\section{Introduction}
A Cartan geometry is a manifold infinitesimally modeled on a
homogeneous space $G/P$, where $G$ is a Lie group and $P$ is a closed
subgroup. It comprises a principal $P$-bundle over the manifold,
equipped with a \emph{Cartan connection} taking values in $\lieg$; see Definition~\ref{defin.cartan.geom} below. Essentially all classical
rigid differential-geometric structures correspond canonically to
Cartan geometries for a suitable homogeneous model, in such a way that the
automorphisms of the structure are the same as the automorphisms of
the Cartan geometry.

Our results concern Cartan geometries, particularly parabolic geometries, having large
automorphism group. This property typically implies further special properties of the geometry, for example,
the existence of solutions to natural overdetermined geometric PDEs. Conversely,
when there are many solutions to such PDEs, one may expect a large automorphism group. Either property implies restrictions on the curvature. The results of this paper are intended to contribute to the concrete understanding of this interplay.

\subsection{Parabolic geometries and tractor bundles}
\label{sec.intro.tractor}

A \emph{parabolic geometry} is a Cartan geometry infinitesimally modeled
on a homogeneous space~$G/P$, where $G$ is a semisimple Lie group and
$P < G$ a parabolic subgroup. This rich family of geometries includes
conformal and projective structures, along with many others. In general, a parabolic
geometry on a manifold $M$ corresponds to a filtration of the tangent
bundle, $TM$, with some additional geometric structures on the
subquotients of the filtration.

Tractor bundles are vector bundles {$\mathcal{V} = \hat{M} \times_\rho {\bf V}$} associated to the Cartan bundle
$\hat{M}$, for $G$-representations $\rho$ restricted to $P$, and equipped
with certain geometrically significant, automorphism-invariant connections.
On the one hand, any tractor bundle is equipped with the tractor connection $\nabla^\omega$, induced by the Cartan connection $\omega$,
and on the other hand with the prolongation connection $\nabla^{\mathcal V}$ of the \emph{first BGG-operator} corresponding to $\mathcal V$, which
in general differs from $\nabla^\omega$ by curvature terms; see Section~\ref{sec_BGG} for details. Parallel sections for these connections correspond to solutions of overdetermined
geometric PDEs
on the manifold; they form a finite-dimensional vector space.

The conformal class of a metric of signature $(p,q)$ on a manifold
$M$, with $p+q = \mbox{dim } M \geq 3$, corresponds to a Cartan geometry
modeled on $G/P$ for $G = \mbox{PO}(p+1,q+1)$ and $P$ the maximal
parabolic subgroup stabilizing an isotropic line in $\BR^{p+1,q+1}$.
The homogeneous space $G/P$ is called the \emph{M\"obius space} of
type $(p,q)$, denoted $\mbox{\bf M\"ob}^{p,q}$, and is diffeomorphic to a two-fold quotient of ${\bf S}^p \times {\bf S}^q$. It carries a
conformally $G$-invariant, conformally flat metric of signature~$(p,q)$. The Cartan
bundle is a principal $P$-subbundle of the second-order frame bundle
of~$M$. The~Cartan connection $\omega$ has been a valuable tool
in recent progress on rigidity of conformal transformations of
semi-Riemannian manifolds (see, for example, \cite{me.notices} for an
overview). Vanishing of the Cartan curvature of this connection over an
open subset $U \subseteq M$ is equivalent to conformal flatness of the
metric on $U$.

A semi-Riemannian metric $g$ on $M^n$, $n \geq 3$, is \emph{Einstein}
if \smash{$\operatorname{Ric} g = \frac{S}{2} g$}, where $\operatorname{Ric}$ is the Ricci
curvature tensor and $S$ is the scalar curvature.
It is \emph{conformally almost-Einstein} if there exists a
smooth function $\sigma$, not constant $0$, and a smooth function $\nu$ such that
\begin{equation}
\label{eqn.conf.einstein}
 \nabla^2 \sigma + \sigma \cdot P = \nu \cdot g,
 \end{equation}
 where $\nabla^2$ is the Hessian, and
 \[
 P = \frac{1}{n-2} \left( \operatorname{Ric} - \frac{S}{2(n-1)} \cdot g \right)
 \]
is the \emph{Schouten tensor}; in other words, $\nabla^2 \sigma +
\sigma \cdot P$ is pure trace. When $\sigma$ is nonvanishing,
then~${\sigma^{-2} \cdot g}$ is an Einstein metric in the conformal class of $g$
(see, e.g., \cite[Proposition~3.6]{curry.gover.tractors}).

Suppose $\hat{M} \rightarrow M$ is the Cartan bundle of a conformal
semi-Riemannian structure. Let ${\bf V} = \BR^{p+1,q+1}$
and $\rho$ the standard representation of ${\rm O}(p+1,q+1)$. The
\emph{(standard) conformal tractor bundle} of such a structure is the associated
vector bundle $\mathcal{V} = \hat{M} \times_\rho {\bf V}$ and it is equipped with the $\omega$-induced tractor connection $\nabla^{\omega}$.
Sections of
$\mathcal{V}$ yield, via projection to a certain line bundle
and a choice of metric in the conformal class,
functions on $M$; functions coming from parallel sections for
$\nabla^{\mathcal{\omega}}$ correspond to
solutions of the conformal almost-Einstein equation. In~this case, the
tractor connection equals the prolongation connection $\nabla^{\mathcal V}$ of the first BGG operator.
See Section~\ref{sec.conf.tractor} below, and the references \cite{cap.gover.tractor.calc, curry.gover.tractors} for more details.

 There are many other significant geometric PDEs on conformal manifolds, and more generally on manifolds equipped with parabolic geometries, whose solutions correspond under prolongation
 to parallel sections of the prolongation connection of a first BGG-operator; on conformal manifolds this applies for instance also to the equations for conformal Killing forms \cite{GoverSilhan, Semmelmann} and for twistor spinors \cite{BFGK}, and for projective and
 c-projective structures to the equation controlling the metrizability of these structures \cite{cemn.cproj, eastwood.matveev.proj.tractor}, among many others.

Infinitesimal automorphisms of geometric structures underlying Cartan geometries also belong to the framework of parallel
sections of natural connections on tractor bundles. Let $\bigl(M, \hat{M}, \omega\bigr)$ be any Cartan geometry modeled on
$(\lieg,P)$. The \emph{adjoint tractor bundle} is the associated vector
bundle $\mathcal{A}$ for ${\bf V} = \lieg$ with the adjoint
representation. A section of $\mathcal{A}$ gives rise via the Cartan
connection to a vector field on $\hat{M}$ which is invariant under the
right-$P$-action, and thus descends to a vector field on $M$.
\v{C}ap observed that for a certain connection on~$\mathcal{A}$,
again related to an intrinsically defined differential operator on $M$, parallel sections correspond to
infinitesimal automorphisms of the Cartan geometry -- that is, vector
fields for which the flow preserves the Cartan connection~\cite{cap.deformations}. These
descend to vector fields on $M$ for which the flow preserves the
geometric structure encoded by the Cartan geometry, called
\emph{infinitesimal automorphisms} or \emph{Killing fields} of the structure.

This article concerns a
generalization of the Gromov--Zimmer embedding theorem, which originally pertains to
infinitesimal automorphisms, to general parallel sections for automorphism-invariant connections of tractor bundles.
We briefly recall the embedding theorem for infinitesimal automorphisms.

\subsection{Gromov--Zimmer embedding theorem}
\label{sec.intro.embedding}

Zimmer's embedding theorem first appeared in \cite{zimmer.lorentz} in the context of a noncompact simple Lie group $H$ acting by automorphisms of a $G$-structure $\omega$ of algebraic type on a compact manifold~$M^n$, such that $G$ also defines a volume form. In this context, the theorem gives a linear injection $\iota\colon \lieh \rightarrow \BR^n$ and an algebraic group $\check{H} < G$ locally isomorphic to $H$ such that $\iota$ intertwines the adjoint representation ${\rm Ad}_\lieh H$ with the representation of $\check{H}$ on $\iota(\lieh) \subseteq \BR^n$; in particular,~$H$~locally embeds in the structure group $G$. Zimmer provided an impressive application for~$\omega$ a~Lorentzian metric, proving that $H$ must be locally isomorphic to ${\rm SL}(2,\BR)$.

Gromov subsequently gave a very general version of the embedding theorem in \cite[Theorem~5.2.A]{gromov.rgs}, which applies to a Lie algebra $\lieh$ of infinitesimal automorphisms of a Gromov-rigid geometric structure $\omega$ of algebraic type on a compact manifold $M$, together with a group $B \leq \Aut(M,\omega)$ preserving a finite measure $\mu$ on $M$ and normalizing $\lieh$. Combined with Gromov's Frobenius theorem \cite[Theorem~1.3.A]{gromov.rgs}, the embedding theorem led to the Gromov--Zimmer centralizer and representation theorems (see \cite{zimmer.gromovrep}).

Our version of the embedding theorem will be a direct generalization of the following version proved by Bader, Frances, and the first author in \cite[Theorems~1.2 and 4.1]{bfm.zimemb} for Cartan geometries.

\begin{thm}%\label{bfm_embedding_thm}
Let $\bigl(M, \hat{M}, \omega\bigr)$ be a Cartan geometry modeled on $(\lieg,P)$ and assume that ${\rm Ad}_{\lieg} P<\operatorname{Aut} \lieg$ is Zariski closed. Let $H\leq\operatorname{Aut}(M,\omega)$ be a Lie subgroup. Let $B \leq H$ preserve a probability measure $\mu$ on $M$. Then for $\mu$-almost-every $x\in M$, for all $\hat x\in\hat M_x$, there exists an algebraic subgroup $\check B<P$ and an algebraic epimorphism
$R\colon \check{B} \rightarrow \bar{B}_d$ such that
\[ (\Ad g) (\iota_{\hat{x}} (X)) = \iota_{\hat{x}} ( R(g). X ) \qquad \forall X \in \mathcal{\lieh},\, g \in \check{B},\]
where $\iota_{\hat{x}} \colon \lieh \rightarrow \lieg$ is the linear injection $X \mapsto \omega_{\hat{x}}(X_{\hat{x}}).$
\end{thm}
\noindent The group $\bar{B}_d$ is the discompact radical of ${\rm Ad}_{\lieh} B \leq \operatorname{Aut} \lieh$, see Definition~\ref{def.discompact.rad} below. Injectivity of $\iota_{\hat{x}}$ is evident from the axioms for the Cartan connection $\omega$, see Definition~\ref{defin.cartan.geom} below.


This embedding theorem says that for $\mu$-almost-every $x \in M$, for all $\hat{x} \in \hat{M}_x$, the injection $\iota_{\hat{x}}$ intertwines the $B$-representation on $\lieh$ with that of an algebraic subgroup $\check{B} < P$ on $\iota_{\hat{x}}(\lieh) \subseteq \lieg$.

\subsection{Results}

Here is our tractor generalization of the embedding theorem. By invariant connection is meant one which is invariant by all automorphisms; in practice, such connections will be defined in a~manner intrinsic to the geometry and invariant for this reason.

\begin{thm}
\label{thm.embedding-0}
 Let $\bigl(M, \hat{M}, \omega\bigr)$ be a Cartan geometry modeled on $(\lieg,P)$. Let $\rho$ be an infinitesimally faithful, completely reducible $G$-representation on $\V$ with $\rho(P)$ algebraic, with associated tractor bundle $\mathcal{V}$, equipped with any invariant connection $\nabla$.
 Let $B \leq \operatorname{Aut}(M,\omega)$ be a Lie subgroup preserving a
 probability measure $\mu$ on $M$.

Let $\mathcal{S}$ denote the $\nabla$-parallel sections of $\cV$.
Then for $\mu$-almost-every $x \in M$, for all $\hat{x} \in \hat{M}_x$, there exists
an algebraic subgroup $\check{B} <P$ with an algebraic epimorphism $R\colon \check{B} \rightarrow \bar{B}_d$ such that
\[ g .\iota_{\hat{x}} (X) = \iota_{\hat{x}} \left( R(g). X \right) \qquad \forall X \in \mathcal{S},\, g \in \check{B}.\]

\end{thm}
\noindent The linear injection $\iota_{\hat{x}}\colon \Gamma(\mathcal{V}) \rightarrow {\bf V}$ is determined by $\hat{x} \in \hat{M}$ as described in Section~\ref{sec.parabolic.tractor}, and $\bar{B}_d$ denotes the discompact radical of the image of $B$ in $\GL(\mathcal{S})$ (see Definition~\ref{def.discompact.rad} below).

\section{Background}

\subsection{Parabolic geometries and tractor connections}

\label{sec.parabolic.tractor}
We will now briefly recall some basics about Cartan geometries, and parabolic geometries in particular.
For a general introduction to Cartan geometries, see \cite{sharpe}.
The comprehensive reference on
parabolic geometries is \cite{cap.slovak.book.vol1}.

\subsubsection{Cartan geometries and parabolic geometries}
%\label{sec.cg.pg}
Suppose $G$ is a Lie group and $P<G$ a closed subgroup. We will always assume that $P$ contains no nontrivial normal subgroup of $G$.
\begin{defin}
 \label{defin.cartan.geom}
 A \emph{Cartan geometry} modeled on $(\g, P)$ on a manifold $M$ is given by a principal $P$-bundle \smash{$\pi\colon \hat{M}\rightarrow M$} equipped with a
 \emph{Cartan connection}, that is, a $P$-equivariant $1$-form ${\omega\in\Omega^1\bigl(\hat{M},\g\bigr)}$ such that
 \begin{itemize}\itemsep=0pt
 \item $\omega_{\hat{x}}\colon T_{\hat{x}} \hat{M}\rightarrow \g$ is an isomorphism for all $\hat{x}\in \hat{M}$;
 \item $\omega(X)= \check{X}$ for any $\check{X}\in \p$,
 where $X_{\hat{x}}= \frac{{\rm d}}{{\rm d}t}\big|_{t=0 } \hat{x}\cdot {\rm e}^{t \check{X}}$ is the vertical vector field on~$\hat{M}$ generated by $\check{X}\in\p$.
 \end{itemize}
 \end{defin}
 The \emph{homogeneous model} of a Cartan geometry as in Definition~\ref{defin.cartan.geom} is the homogeneous space~$G/P$ equipped with the natural projection $G\rightarrow G/P$ and
 the Maurer--Cartan form ${\omega_G\in\Omega^1(G,\g)}$ of $G$, which encodes the natural trivialization $TG\cong G\times\g$ by left-invariant vector fields.
 The \emph{curvature of a Cartan geometry} is given by the following semibasic, $\g$-valued, $P$-equivariant $2$-form on $\hat{M}$:
 \begin{equation*}
 K(\xi,\eta)={\rm d}\omega(\xi,\eta)+[\omega(\xi), \omega(\eta)] \qquad\textrm{for } \xi,\eta\in\mathfrak X(\hat{M}).
 \end{equation*}
 The significance of the curvature is that it vanishes if and only if the Cartan geometry is locally isomorphic to its homogeneous model.

 Given a Cartan geometry, any representation $\rho\colon P\rightarrow\textrm{GL}({\bf V})$ of $P$ induces
 a vector bundle by forming the associated bundle to the Cartan bundle:
 \[
 \mathcal V =\hat{M} \times_\rho {\bf V}.
 \]
 Recall that the space of smooth sections $\Gamma(\mathcal V)$ of $\mathcal V$ can be identified with the space \smash{$C^\infty\bigl(\hat{M}, {\bf V}\bigr)^P$} of smooth, $P$-equivariant functions $f\colon \hat{M}\rightarrow {\bf V}$, that is,
 \[f(\hat{x}\cdot p)= {\rho\bigl(p^{-1}\bigr). f(\hat{x})}.\]
 In this way, each $\hat{x} \in \hat{M}$ gives a linear function $\iota_{\hat{x}}\colon \Gamma(\mathcal{V}) \rightarrow {\bf V}$ by evaluation.

 Sections of the adjoint bundle $\mathcal{A} = \hat{M} \times_P \lieg$ from Section~\ref{sec.intro.tractor} correspond via $\omega$ to right-$P$-invariant vector fields on $\hat{M}$, which descend to vector fields on $M$.
The Cartan connection $\omega$ induces isomorphisms of vector bundles
 \begin{equation}\label{iso_tangent_bundle}
 TM\cong \hat{M} \times_P \g/\p\qquad\textrm{and}\qquad T^*M\cong \hat{M} \times_P (\g/\p)^*.
 \end{equation}
 These identifications can easily be derived from the axioms for $\omega$ (see also \cite[Theorem~5.3.15]{sharpe}).
 Because the curvature $K$ is semibasic and $P$-equivariant, it
 descends to a $2$-form on $M$ with values in $\mathcal{A}$, which, via \eqref{iso_tangent_bundle}, can be identified with a $P$-equivariant \emph{Cartan curvature function}
 \begin{equation*}
 \kappa\colon\ \hat{M} \rightarrow\Lambda^2 (\g/\p)^*\otimes \g.%\cong \Lambda^2 \p_+\otimes \g.
 \end{equation*}

Many important geometric structures admit equivalent descriptions as Cartan geometries, in particular semi-Riemannian metrics. A broad and interesting class of examples are the geometric structures underlying \emph{parabolic geometries}, which are Cartan geometries of type $(\g,P)$, where $P<G$ is a parabolic subgroup of a semisimple Lie group.
These correspond to certain bracket-generating distributions on $M$ endowed with a geometric structure defined by $G_0$, the reductive complement in a Levi decomposition of $P \cong G_0 \ltimes P_+$. These structures correspond canonically to \emph{normal} parabolic Cartan geometries. See \cite{cap.schichl.equiv,morimoto,tanaka} for the precise general statements, including the notions of normal and \emph{regular}.
In the table below, we list some of the most prominent examples of geometric structures admitting an equivalent description as regular normal parabolic geometries of type $(\g,P)$. For more examples, see \cite{cap.slovak.book.vol1}. A concise reference on the construction
of the Cartan connection in conformal or projective geometry is \cite[Section IV]{kobayashi.transf}. All parabolic geometries will be assumed to be regular and normal throughout the sequel.

\begin{table}[h]\centering
\small\renewcommand{\arraystretch}{1.15}
\begin{tabular}{ |l|c|l| }
\hline
 Geometric structure & $\g$ & $P$ \\
\hline\hline
Projective structures on $n$-manifolds & $\mathfrak{sl}(n+1,\BR)$& stabilizer of line in $\BR^{n+1}$\\
\hline
Almost c-projective structures &$\mathfrak{sl}(n+1,\C)$ &stabilizer of complex \\
on $2n$-manifolds, $n\geq 2$ &&line in $\C^{n+1}$\\
\hline
Conformal $n$-manifolds of &$\mathfrak{so}(p+1,q+1)$ & stabilizer of an isotropic \\
signature $(p,q)$, $n=p+q\geq 3$ & &line in $\BR^{p+1,q+1}$ \\
 \hline
 Almost quaternionic $4n$-manifolds& $\mathfrak{sl}(n+1,{\bf H})$ & stabilizer of a quaternionic \\
 &&line in ${\bf H}^{n+1}$ \\
 \hline
 Non-degenerate partially integrable &$\mathfrak{su}(p+1,q+1)$& stabilizer of an isotropic \\
 CR-structures of hypersurface type & &complex line in Hermitian \\
 of signature $(p,q)$ &&vector space $\C^{p+1,q+1}$\\
 \hline
 Generalized path geometries&$\mathfrak{sl}(n+2,\R)$& stabilizer in $\BR^{n+2}$ of a flag \\
 on $2n+1$-manifolds, $n\geq 2$ && of line inside a plane\\
 \hline
 (2,3,5)-distributions on $5$-manifolds &$G_2$&stabilizer of highest weight\\
 && line in the $7$-dimensional \\
 &&fundamental representation\\
 \hline
\end{tabular}
\end{table}


On the Lie algebra level, a Levi decomposition will be denoted $\p \cong \g_0 \ltimes \p_+$. There is a $G_0$-invariant subspace of $\g$ complementary to $\p$, which is denoted $\g_-$; it is of course isomorphic as a~$G_0$-module to $\g/\p$. The nilradical $\p_+ \lhd \p$ is moreover isomorphic as a $P$-module to $(\g/\p)^*$. Any representation $\bf V$ of $\p$ on which the center of $\g_0$ acts diagonalizably
has
 a $P$-invariant filtration
\begin{equation}\label{filtration_rep}
{\bf V}={\bf V}^{-\ell}\supset {\bf V}^{- \ell+1}\supset \cdots \supset {\bf V}^0\supset {\bf V}^{1}\supset \cdots \supset {\bf V}^{m} \supset 0
\end{equation}
such that each quotient ${\bf V}^{i}/{\bf V}^{i+1}$ is a completely reducible $\p$-representation, which implies that~$\p_+$ acts trivially on ${\bf V}^{i}/{\bf V}^{i+1}$.
In particular, the $\p$-representation $\g$ admits such a filtration
 \begin{equation*}
 \g=\g^{-k}\supset\dots \supset \g^0\supset\dots \supset \g^k,
 \end{equation*}
 where $\g^0=\p$ and $\g^1=\p_+$, and so does the isotropy representation
 \begin{equation}\label{filtr_g/p}
 \g/\p=\g^{-k}/\p\supset\dots \supset \g^{-1}/\p.
\end{equation}
Given two filtered vector spaces $\bf V$ and $\bf W$, there is an induced filtration by \emph{homogeneity} on the linear maps $\operatorname{Hom}({\bf V},{\bf W})={\bf V}^*\otimes {\bf W}$,
in which $\phi\in\operatorname{Hom}({\bf V},{\bf W})^i$ if $\phi({\bf V}^j)\subset {\bf W}^{i+j}$ for all $j$.

\subsubsection{Tractor connections}
\label{sec.tractor.connxns}
For a Cartan geometry $\bigl(M,\hat{M}, \omega\bigr)$ of any type $(\g,P)$, the natural vector bundles associated to $G$-representations play a distinguished role.
The Cartan connection $\omega$ can be canonically extended to a principal connection on the principal $G$-bundle $\tilde{M}= \hat{M} \times_PG$ and as such it induces a linear connection
$\nabla^\omega$ on all vector bundles of the form
\[\mathcal V=\tilde{M} \times_G{\bf V}= \hat{M} \times_P{\bf V}.\]
\begin{defin} Given a Cartan geometry,
the vector bundles associated to $G$-representations are called \emph{tractor bundles} and $\nabla^\omega$ the \emph{tractor connection}.
\end{defin}

The adjoint tractor bundle comes with a natural projection $\Pi\colon \mathcal A\rightarrow TM$ to the tangent bundle induced by the projection $\g \rightarrow \g/\p$,
and an algebraic action on any other tractor bundle $\mathcal A\times \mathcal V\rightarrow\mathcal V$, denoted by $(s,t)\mapsto s\bullet t$, induced by the action of~$\g$ on ${\bf V}$.

For any vector bundle $\mathcal W$ associated to any $P$-representation ${\bf W}$, there is a natural differential operator, called the \emph{fundamental derivative}, denoted by
 \begin{align*}
 D\colon\ \Gamma(\mathcal A)\times\Gamma(\mathcal W) \rightarrow \Gamma(\mathcal W),\qquad
(s,t) \mapsto D_st.
 \end{align*}
Viewing $t$ in \smash{$C^\infty\bigl(\hat{M}, {\bf W}\bigr)^P$} and $s$ in \smash{$C^\infty\bigl(\hat{M},\g\bigr)^P$}, the section \smash{$D_st \in C^\infty\bigl(\hat{M},{\bf W}\bigr)^P$} evaluates at $\hat{x} \in \hat{M}$ to the derivative of $t$ in the direction of $\omega_{\hat{x}}^{-1}(s(\hat{x}))$. The tractor connection on a tractor bundle $\mathcal V$ can then be expressed in terms of the fundamental derivative as follows:
\begin{equation}
\label{eqn.def.tractor.connxn}
\nabla_{\Pi(s)}^\omega t=D_st+s\bullet t\qquad \textrm{for}\ s\in \Gamma(\mathcal A),\, t\in\Gamma(\mathcal V).
\end{equation}
It is a routine verification that this expression defines a connection; in particular, at any point, the right-hand side only depends on $\Pi(s)$ at that point and on $t$ in a neighborhood. If $t$ is a~parallel section, it is the unique one with a given value $t(x)$ at any $x \in M$; it follows that $\iota_{\hat{x}}$ is an injective linear map to ${\bf V}$ when restricted to parallel sections in $\Gamma(\mathcal{V})$, for any $\hat{x} \in \hat{M}$.

As discussed in the introduction, parallel sections of tractor connections typically correspond to solutions of interesting equations for the geometric structure on $M$ described by the Cartan geometry. More precisely, in many contexts, certain deformations of the tractor connections by curvature terms yield invariant connections on the tractor bundles for which the parallel sections correspond to the solutions to a particular geometric PDE of interest.
 In the case of parabolic geometries, there is a powerful framework for constructing such connections, which we now briefly present.
\subsubsection{First BGG-operators of parabolic geometries}\label{sec_BGG}
Suppose $\bigl(M,\hat{M},\omega\bigr)$ is a regular normal parabolic geometry of type $(\g,P)$ with $\dim M=n$. It was shown in \cite{CD.BGG, css.bgg.seqs} that any $G$-representation ${\bf V}$ induces a \emph{BGG-sequence}
of intrinsically defined differential operators of the form
\begin{equation}\label{bgg.seq}
 \Gamma(\mathcal H_0)\stackrel{\mathcal D^{\mathcal V}_0\ }{\rightarrow}\Gamma(\mathcal H_1)\overset{\mathcal D^{\mathcal V}_1\ }{\rightarrow}\cdots
 \overset{\mathcal D^{\mathcal V}_{n-1}}{\rightarrow}\Gamma(\mathcal H_{n}),
 \end{equation}
 where \[\mathcal H_k = \mathcal H_k(M,\mathcal V)=\hat{M}\times_P H_k(\p_+,{\bf V})\] is the bundle associated to the $k$-th homology $H_k(\p_+,{\bf V})$ of $\p_+$ with coefficients in ${\bf V}$. By \cite{kostant.harmonic.curv},
 the bundles $\mathcal H_k$ are weighted tensor bundles on $M$. In the case of the homogeneous model $M \cong G/P$, the BGG-sequence is a complex, which is dual to the \emph{BGG-resolution} of ${\bf V}$ by generalized Verma modules, whence the terminology.

 The first operator $\mathcal D^{\mathcal V} =\mathcal D^{\mathcal V}_0\colon \Gamma(\mathcal H_0)\rightarrow \Gamma(\mathcal H_1)$ in \eqref{bgg.seq} always defines an overdetermined system of PDEs, the
 domain of which is easily computed to be $\mathcal H_0=\mathcal V/\mathcal V^+,$ where $\mathcal V^+ \subset \mathcal V$ is the subbundle corresponding to the $P$-submodule $\p_+{\bf V} \subset {\bf V}$.
 We denote by $\Pi^\mathcal V\colon \mathcal V\rightarrow \mathcal V/\mathcal V^+=\mathcal H_0$ the natural projection. Then, \cite{hsss.prol.connxn, neuss.prolong} shows that the projection $\Pi^\mathcal V$ restricts to a bijection between
 sections of $\mathcal V$ that are parallel for a particular linear connection on $\cV$ and solutions to~\smash{$\mathcal D^{\mathcal V}(\sigma)=0$}.

 The linear connection on $\mathcal V$, as constructed in \cite{hsss.prol.connxn}, the parallel sections of which correspond to solutions of $\mathcal D^{\mathcal V}(\sigma)=0$, is called the \emph{prolongation connection} of $\mathcal D^{\mathcal V}$ and will be denoted by~$\nabla^\mathcal V$. In general, $\nabla^\mathcal V$ differs from the tractor connection $\nabla^\omega$ by tensors built from the curvature of~$\nabla^\omega$. As above, the linear map $\iota_{\hat{x}}$ is injective when restricted to parallel sections for $\nabla^{\mathcal{V}}$ for all~${\hat{x} \in \hat{M}}$.

\subsubsection{Automorphisms acting on tractor bundles}
\label{sec.aut.tractor}

The \emph{automorphisms} of a Cartan geometry $\bigl(M,\hat{M},\omega\bigr)$ are the diffeomorphisms of $M$ that lift to principal bundle automorphisms of $\hat{M}$ which preserve $\omega$. For $f \in \operatorname{Aut}(M,\omega)$, we will also denote by $f$ its lift to $\hat{M}$, satisfying $f^* \omega = \omega$.
The automorphisms of a Cartan geometry naturally preserve the underlying geometric structure; for example, the automorphisms of the Cartan geometry associated to a conformal semi-Riemannian structure are precisely the conformal transformations.

Because it preserves the framing by $\omega$, the action of $\Aut(M,\omega)$ on $\hat{M}$ is free and proper. If $f \in \Aut(M,\omega)$ stabilizes $x \in M$, then for any $\hat{x} \in \hat{M}_x$, there is a unique $p \in P$ such that {$f.\hat{x} = \hat{x}.p$}, called the \emph{isotropy of $f$ at $x$ with respect to $\hat{x}$}.

For any $P$-module ${\bf W}$, the group $\Aut(M,\omega)$ acts on the associated bundle $\mathcal{W} = \hat{M} \times_P {\bf W}$.
The action of $f \in \Aut(M,\omega)$ on a section \smash{$X \in \Gamma(\cS)\cong C^{\infty}\bigl(\hat{M}, {\bf W}\bigr)^P$} is simply by
$f^*X = X \circ f^{-1}$.

Assume now that $\bigl(M,\hat{M},\omega\bigr)$ is a regular parabolic geometry modeled on $(\g, P)$, and that $\rho\colon G\rightarrow \textrm{GL}(\V)$ is a $G$-representation.
Let $\cV=B\times_P\V$ be the corresponding tractor bundle with tractor connection $\nabla^{\omega}$. Let $\nabla^{\cV}$ be the prolongation connection of the
first BGG-operator~$D^{\cV}$ corresponding to $\cV$.
\begin{prop} Any $f \in\emph{Aut}(M, \omega)$ preserves both $\nabla^\omega$ and $\nabla^{\cV}$, that is, $f^* \nabla^\omega = \nabla^\omega$ and~${f^*\nabla^{\cV}=\nabla^{\cV}}$.
\end{prop}

\begin{proof}
Since $f^*\omega=\omega$ for any $f \in \Aut(M,\omega)$, the fundamental derivative is invariant:
\[D_{f^* X} f^* Y = f^* (D_X Y) \qquad \forall X \in \Gamma({\mathcal{A}}), \, Y \in \Gamma(\mathcal{V}).\]
The algebraic action of $\Gamma(\mathcal{A})$ on $\Gamma(\mathcal{V})$ is evidently invariant by any principal bundle automorphism.
Thus it follows from (\ref{eqn.def.tractor.connxn}) that $f^*\nabla^\omega=\nabla^{\omega}$.

Set $\Lambda=\nabla^{\cV}-\nabla^{\omega}\in\Omega^1(M, \operatorname{End}(\cV))$ and denote by $R^{\cV}\in\Omega^2(M, \operatorname{End}(\cV))$ the curvature of~$\nabla^{\cV}$. The differential computing the Lie algebra homology of $\p_+$ with values in $\V$ induces via the identification
\[T^*M \cong \hat{M} \times_P (\g / \p)^* \cong \hat{M} \times_P \p^+\]
 a bundle map
$\partial^*_{\cV}\colon \Omega^\ell(M, \cV)\rightarrow \Omega^{\ell-1}(M, \cV)$. The connection $\nabla^{\cV}$ is characterized by the two properties \cite{hsss.prol.connxn}:
\begin{itemize}\itemsep=0pt
\item[(1)] $\Lambda\in\Omega^1(M, \operatorname{End}(\cV))^1\cap \Gamma(\operatorname{Im}(\partial^*_{\cV}\otimes {\rm Id}_{\cV^*}))$,
\item[(2)] $R^{{\cV}}\in\Gamma(\ker(\partial^*_{\cV}\otimes {\rm Id}_{\cV^*})).$
\end{itemize}
The vector bundle $(T^*M\otimes\operatorname{End}(\cV))^1$ is the subbundle of
\[
T^*M\otimes\operatorname{End}(\cV)\cong \hat M\times_P\operatorname{Hom}(\g/\p,\operatorname{End}({\bf V}))
\]
 corresponding to the first filtration component, that is, the homogeneity-one component, of the filtered $P$-representation
$\operatorname{Hom}(\g/\p,\operatorname{End}({\bf V}))$, where the filtration on that space is induced by~\eqref{filtration_rep} and~\eqref{filtr_g/p}; as it is an associated bundle to the Cartan bundle, it is acted on by $\Aut(M,\omega)$. The subbundle $\operatorname{Im}(\partial^*_{\cV}\otimes {\rm Id}_{\cV^*})$ is similarly $\Aut(M,\omega)$-invariant, so any $f \in \Aut(M,\omega)$ preserves their intersection, thus $f^*\Lambda$
is a section of this intersection. As $f^*\Lambda = f^*\nabla^\cV-\nabla^{\omega}$, the pullback~$f^* \nabla^\cV$ satisfies (1).
Similarly, the curvature of $f^*\nabla^\cV$, namely $f^*R^{{\cV}}$, is a section of the associated vector bundle $\ker(\partial^*_{\cV}\otimes {\rm Id}_{\cV^*})$, so $f^* \nabla^\cV$ satisfies (2) as well.
We conclude $f^*\nabla^{\cV}=\nabla^{\cV}$ thanks to this characterization.
\end{proof}

\subsubsection{The conformal tractor connection}
\label{sec.conf.tractor}

Let $(M,g)$ be a semi-Riemannian manifold of dimension $n \geq 3$ and let $\bigl(M, \hat{M}, \omega\bigr)$ and $\mathcal{V}$ be the corresponding parabolic Cartan geometry with model $(\g, P)$ and conformal tractor bundle, respectively, as in Section~\ref{sec.intro.tractor}. Recall that $P$ is the stabilizer of an isotropic line ${\bf V}_{1} \subset {\bf V} \cong \BR^{p+1,q+1}$. Setting ${\bf V}^{1} = {\bf V}_{1}$ and {${\bf V}^0 = {\bf V}_{1}^\perp$}, there is a $P$-invariant filtration
${{\bf V}^{1} \subset {\bf V}^0 \subset {\bf V}^{-1} = {\bf V}}$, cf.\ \eqref{filtration_rep}, giving rise to an intrinsic filtration $ \mathcal{V}^{1} \subset \mathcal{V}^0 \subset \mathcal{V}$.

A Levi complement $G_0 < P$ preserves a splitting ${\bf V} = {\bf V}_{1} \oplus {\bf V}_0 \oplus {\bf V}_{-1}$.
A choice of metric $g_0 \in [g]$ gives a corresponding decomposition
\begin{equation}
\label{eqn.conf.splitting}
\mathcal{V} = \mathcal{V}_{1} \oplus \mathcal{V}_0 \oplus \mathcal{V}_{-1}
\end{equation}
in which $\mathcal{V}_0 \cong TM$ and $\mathcal{V}_{-1}$ is identified with the line bundle of conformal scales.
References for the formulas in this section are \cite[equations~(22) and~(32)]{curry.gover.tractors}.
With respect to a decomposition~(\ref{eqn.conf.splitting}), the standard conformal tractor connection $\nabla^\omega$ is
\begin{equation}
\label{eqn.conf.connxn}
 \nabla^\omega \colon\ \begin{pmatrix} \nu \\ X \\ \sigma \end{pmatrix} \mapsto \begin{pmatrix} {\rm d} \nu - P \llcorner X \\ \nabla X + \nu \cdot {\rm Id} + \sigma P^\sharp \\ {\rm d} \sigma - X^* \end{pmatrix},
 \end{equation}
where $P^\sharp$ is the symmetric endomorphism corresponding to the Schouten tensor $P$, and all duals are with respect to $g_0$. Changing the metric $g_0 \in [g]$ corresponds to applying an element of $P \backslash G_0$ to the decomposition of tractor sections, and this transformation commutes with the tractor covariant derivative.
It follows that the intrinsic $(p+1,q+1)$-metric on $\mathcal{V}$ is $\nabla^\omega$-parallel.

Continuing with $g_0 \in [g]$, the $\sigma$-component of a $\nabla^\omega$-parallel section satisfies
\[ \nabla \mbox{grad } \sigma + \sigma P^\sharp = - \nu \cdot {\rm Id},\]
which is equivalent to (\ref{eqn.conf.einstein}) (with the sign of $\nu$ switched). Parallel sections have constant self-scalar products with respect to the tractor metric and in particular have a well-defined causal~type.

The curvature of $\nabla^\omega$ has the form, in the decomposition (\ref{eqn.conf.splitting}),
\begin{equation}
\label{eqn.conf.curv}
R^\omega(Y,Z) \colon\ \begin{pmatrix} \nu \\ X \\ \sigma \end{pmatrix} \mapsto \begin{pmatrix} - C(Y,Z)X \\ \sigma C(Y,Z)^* + W(Y,Z) X \\ 0 \end{pmatrix},
\end{equation}
where $W$ is the Weyl curvature and $C$ the Cotton tensor of $g_0$:
\[C(Y,Z)X = \frac{1}{2} \left( (\nabla_Y P)(Z,X) - (\nabla_Z P)(Y,X) \right).\]

\subsection{Borel density theorem}

The proof of the embedding theorem rests on the Borel density
theorem.
The term
\emph{algebraic} here means real-algebraic. A Zariski closure in this setting comprises the real points of the Zariski closure over ${\bf C}$.

 We use a version of the Borel density theorem which can be found in \cite[Corollary~3.2]{bfm.zimemb}, and which is due in this generality to Shalom
\cite[Theorem~3.11]{shalom.discompact} (see also \cite[Theorem~2.6]{dani.borel.density}).
The setting is a locally compact group $S$ acting on an algebraic variety $\mathbb{W}$ via a
continuous homomorphism $\rho \colon S \rightarrow \mbox{Aut } \mathbb{W}$. We
will denote by
$\bar{S}$ the Zariski closure of $\rho(S)$ in $\mbox{Aut }\mathbb{W}$.

\begin{defin}
 \label{def.discompact.rad}
 Let $G$ be an algebraic group.
The \emph{discompact radical} of $G$ is the maximal algebraic subgroup
of $G$ with no proper cocompact, algebraic, normal subgroups, denoted $G_d$.
 \end{defin}

 The discompact radical can be obtained as the minimal algebraic
 subgroup of $G$ containing all algebraic subgroups with no non-trivial compact
 algebraic quotients, which exists by the Noetherian property of
 algebraic subgroups. See \cite{shalom.discompact} for more details.

\begin{thm}[{\cite[Corollary 3.2]{bfm.zimemb}} of the Borel density theorem]
 \label{thm.borel.density}
Let $S$ be a locally compact group and $\mathbb{W}$ an algebraic variety, with
an action $\rho\colon S \rightarrow \mbox{Aut } \mathbb{W}$.
Suppose $S$ acts by homeomorphisms on a topological space $M$
preserving a finite Borel measure $\mu$. Assume $\phi\colon M \rightarrow
\mathbb{W}$ is an $S$-equivariant measurable map. Then $\phi(x)$ is fixed
by $\bar{S}_d$ for $\mu$-almost-every $x \in M$.
 \end{thm}

\section{Embedding theorem}

The following theorem is the tractor generalization of Theorem 1.2 of
\cite{bfm.zimemb} (see also Theorem~4.1 there), restated here as in the introduction with an additional technical point in the conclusions. The proof will be a straightforward adaptation of the one in~\cite{bfm.zimemb}.


\begin{thm}
\label{thm.embedding}
Let
$\bigl(M, \hat{M},\omega\bigr)$ be a Cartan geometry modeled on $(\lieg,P)$ and $B \leq \Aut(M,\omega)$ a~Lie subgroup preserving a Borel probability measure $\mu$.
Let $\rho$ be a fully reducible representation of~$G$ on $\V$ with $\rho(P)$ Zariski closed, and $\mathcal{V} = \hat{M}
 \times_\rho \V$. Let $\mathcal{S}$ be the space of parallel sections for any $\Aut(M,\omega)$-invariant connection $\nabla$ on $\mathcal{V}$.

 Then for $\mu$-almost-every $x \in M$, for all $\hat{x} \in \hat{M}_x$,
there is an algebraic subgroup $\check{B}_{\hat{x}} \leq \rho(P)$ and an algebraic epimorphism $R\colon \check{B}_{\hat{x}} \rightarrow \bar{B}_d$ such that
\[ g. \iota_{\hat{x}}(X) = \iota_{\hat{x}} {( R(g).X)} \qquad \forall X \in \mathcal{S}, \ g \in \check{B}_{\hat{x}},\]
moreover, the isotropy of $B$ at $x$ with respect to $\hat{x}$ is contained in $\check{B}_{\hat{x}}$.
\end{thm}


Recall that $\bar {B}_d$ denotes the discompact radical of the Zariski closure of
 the image of $B$ in~$\GL(\cS)$.


 \begin{proof}
Denote by $\psi\colon \Aut(M,\omega) \rightarrow \textrm{GL}(\cS)$ the representation corresponding to the restriction to $\mathcal{S}$ of the action on $\Gamma(\mathcal{V})$ as in Section~\ref{sec.aut.tractor}.
Let
\begin{align*}
\iota \colon \ \hat{M} \rightarrow \textrm{Mon}(\cS, \V), \qquad
 \hat{x} \mapsto \iota_{\hat{x}},
\end{align*}
where $ \textrm{Mon}(\cS, \V)$ is the variety of all linear injections
$\cS \rightarrow \V$. Each $g \in \Aut(M,\omega)$ acts on $\textrm{Mon}(\cS,\V)$ by pre-composition with $\psi(g)$, an algebraic automorphism, while each $p \in P$ acts by post-composition with $\rho(p)$, also an algebraic automorphism. These two actions commute.

If $g \in \Aut(M,\omega)$ stabilizes $x$ and has isotropy with respect to $\hat{x}$ equal to $p \in P$, then
\begin{equation}
\label{eqn.isotropy.in.bcheck}
\iota_{\hat{x}}(g^*X) = {p.\iota_{\hat{x}}(X)}.
\end{equation}


Let $\mathbb{U} = \mbox{Mon}(\cS, \V)/P$ and $\bar{\iota}\colon M \rightarrow \mathbb{U}$ the induced map, which is also $\Aut(M,\omega)$-equivariant.
Note that if $g \in \Aut(M,\omega)$ stabilizes $x$, then $\psi(g)$ stablizes $\bar{\iota}(x)$.
Although $\mathbb{U}$ is not generally an
algebraic variety, the Rosenlicht stratification
\cite{rosenlicht.brasil} (see also \cite[Section~2.2]{gromov.rgs}) yields a
decomposition
\[\mathbb{U} = \mathbb{U}_0 \cup \cdots \cup \mathbb{U}_r\]
in which
each $\mathbb{U}_i$ is a variety on which $\Aut(M,\omega)$ acts via $\psi$ by automorphisms
(see \cite[Theorem~4.1]{bfm.zimemb} for more details).
The sets $M_i = \bar{\iota}^{\,-1}(\mathbb{U}_i)$ are
$B$-invariant and measurable for each $i=0, \dots, r$. For each $i$ such that $\mu_i = \left. \mu \right|_{M_i}$ is not zero, it is a nontrivial finite,
$B$-invariant measure.


By Theorem~\ref{thm.borel.density}, for $\mu_i$-almost-every $x\in M_i$ the point $\bar{\iota}(x)$ is fixed by
$\bar B_d$.
Let $F_i$ comprise the $\bar{B}_d$-fixed points in
$\mathbb{U}_i$, and set $\Omega_i = (\bar{\iota} \circ \pi)^{-1}(F_i)$, a $P$-invariant subset
of $\hat{M}$ with ${\mu_i(M_i \backslash \pi(\Omega_i)) = 0}$.
For any $\hat{x} \in \Omega_i$, the orbits
\[ \iota(\hat{x}).\bar B_d \subset P.\iota(\hat{x}).\]
If we now set
\[\check{B}_{\hat{x}}=\bigl\{p\in \rho(P)\mid p.\iota(\hat{x})=\iota(\hat{x}).g \textrm{ for some } g\in\bar B_d\bigr\},\]
then $\check{B}_{\hat{x}}$ is an algebraic subgroup of $\rho(P)$ satisfying the conclusions of
the theorem for $\hat{x}$; in particular, it contains the isotropy of $B$ with respect to $\hat{x}$ by (\ref{eqn.isotropy.in.bcheck}). Let $\Omega = \cup_i \Omega_i$, a full measure set for $\mu$ comprising points that all satisfy the conclusions of the theorem.
\end{proof}

\begin{rem}
The conclusions of the theorem clearly hold when $\mathcal{S}$ is replaced by any $B$-invariant subspace of $\mathcal{S}$.
\end{rem}

For future use, we make the following general definition.

\begin{defin}
\label{def.rep.embedding}
Let $\psi\colon B \rightarrow \GL({\W})$ and $\rho \colon P \rightarrow \GL({\V})$ be two representations, with $\rho(P)$ algebraic. An \emph{embedding of $({\W}, \psi(B))$ into $({\V}, \rho(P))$} comprises a linear injection $\iota\colon {\W} \rightarrow {\V}$ together with an algebraic subgroup $\check{B} \leq \rho(P)$ and an algebraic epimorphism $R\colon \check{B} \rightarrow \bar{B}_d$ such~that
\[ g. \iota(X) = \iota( R(g).X) \qquad \forall X \in {\bf W}, \, g \in \check{B}.\]
\end{defin}

As before, $\bar{B}_d$ denotes the discompact radical of the Zariski closure of $\psi(B)$.
Note that whenever there is a Lie group monomorphism $H \rightarrow G$ for which $\Ad G$ is algebraic, then there is an embedding of $(\lieh, \Ad H)$ into $(\g, \Ad G)$.

\begin{rem}
Under the hypotheses of Theorem~\ref{thm.embedding}, let
$K \colon \hat{M} \rightarrow \mathbb{W}$ be an $\Aut(M,\omega)$-invariant, $P$-equivariant map to a variety $\mathbb{W}$ on
 which $P$ acts algebraically.
 In the proof of the theorem, we can take
\[\mathbb{U} =\textrm{Mon}(\cS, \V)\times \mathbb{W}\]
with $\Aut(M,\omega)$ acting by precomposition via $\psi$ on the first factor and trivially on the second factor, commuting with the $P$-action on the product.
Considering the map $\iota \times K \colon \hat{M}\rightarrow \mathbb{U}$ given
by $\hat{x} \mapsto (\iota_{\hat{x}}, K(\hat{x}))$, the same proof with the Borel density theorem as above yields, for almost-every $x \in M$, not only the embedding of $(\mathcal{S}, \psi(B))$ into $({\bf V}, \rho(P))$ but also that $K(\hat{x})$ is fixed by $\check{B}_{\hat{x}}$; the conclusion that $\check{B}_{\hat{x}}$ contains the isotropy of $B$ at $x$ with respect to $\hat{x}$ moreover remains true.
\end{rem}

A useful choice of $\mathbb{W}$ in the remark above is, for $m \in {\bf N}$, the $P$-module in which the total
order-$m$ fundamental derivative of the Cartan curvature
\[D^m \kappa\colon \ \hat{x} \mapsto \bigl( \kappa(\hat{x}), D \kappa(\hat{x}), \dots, D^{(m)} \kappa(\hat{x}) \bigr)\]
takes its
values (see Sections~\ref{sec.tractor.connxns} and \ref{sec.aut.tractor}). Fixed points of $P$ in this module can give rise to local automorphisms
by the Frobenius theorem for Cartan geometries. More precisely, for~$\bigl(M,\hat{M}, \omega\bigr)$ real-analytic and $M$ closed, there is $m \geq 0$ such that for any $\hat{x} \in \hat{M}$, any $p \in P$ fixing $D^m \kappa(\hat{x})$ is the isotropy at $x = \pi(\hat{x})$ with respect to $\hat{x}$ of a unique $g \in \Aut^{\rm loc}(M,\omega)$ fixing~$x$ (where uniqueness applies to any two such local automorphisms defined in the same neighborhood of~$x$) \cite{me.frobenius}, see also \cite{pecastaing.frobenius} for a smooth version. The following corollary of the embedding theorem is well known in the setting of infinitesimal automorphisms (see \cite[Theorem~5.3]{me.frobenius}).

\begin{cor}
\label{cor.emb.frob}
Let
$\bigl(M, \hat{M},\omega\bigr)$ be a real-analytic Cartan geometry modeled on $(\lieg,P)$ and $B \leq \Aut(M,\omega)$ a Lie subgroup preserving a Borel probability measure $\mu$.
Let $\rho$ be a fully reducible, real-analytic representation of $G$ on $\V$ with $\rho(P)$ Zariski closed, and $\mathcal{V} = \hat{M}
 \times_\rho \V$. Let $\mathcal{S}$ be the space of parallel sections for any $\Aut(M,\omega)$-invariant connection $\nabla$ on $\mathcal{V}$.

 Then for $\mu$-almost-every $x \in M$, for all $\hat{x} \in \hat{M}_x$, in addition to the conclusions of Theorem~$\ref{thm.embedding}$, every $g \in \check{B}_{\hat{x}}$ is the isotropy with respect to $\hat{x}$ of an element of $\Aut^{{\rm loc}}(M,\omega)$ fixing~$x$.
\end{cor}

Parallel tractor sections, like infinitesimal automorphisms, correspond to solutions of overdetermined PDEs, and a Frobenius theorem in terms of fundamental derivatives of the tractor curvature is possible. We have worked out such a result, however, we have not seen how to profitably combine it with the output of the embedding theorem, so it is not included in the present paper.

\begin{rem}
\label{rem.conjugate.embedding}
Let $\hat{x}$ be as in the conclusion of Corollary \ref{cor.emb.frob}. The statement says the conclusions hold for $\hat{x}.g$ for all $g \in P$, as well. In this case, the isotropy of $B$ with respect to $\hat{x}.g$ is that with respect to $\hat{x}$, conjugated by $g^{-1}$. We also have
\[D^m \kappa(\hat{x}.g) = g^{-1}. D^m \kappa(\hat{x}) \qquad \mbox{and} \qquad \check{B}_{\hat{x}.g} = g^{-1} \check{B}_{\hat{x}} g.\]
The group $\check{B}_{\hat{x}}$ will thus be conjugated in $P$ as needed below, corresponding to choosing a~different~${\hat{x} \in \hat{M}_x}$, without altering the properties ensured by the embedding theorem.
\end{rem}


\begin{thebibliography}{99}
\footnotesize\itemsep=-0.5pt

\bibitem{alekseevski.selfsim}
Alekseevski D., Self-similar {L}orentzian manifolds,
 \href{https://doi.org/10.1007/BF00054491}{\textit{Ann. Global Anal. Geom.}}
 \textbf{3} (1985), 59--84.

\bibitem{amores.killing}
Amores A.M., Vector fields of a~finite type {$G$}-structure,
 \href{https://doi.org/10.4310/jdg/1214434847}{\textit{J.~Differential
 Geometry}} \textbf{14} (1979), 1--6.

\bibitem{bfm.zimemb}
Bader U., Frances C., Melnick K., An~embedding theorem for automorphism groups
 of {C}artan geometries,
 \href{https://doi.org/10.1007/s00039-009-0002-x}{\textit{Geom. Funct. Anal.}}
 \textbf{19} (2009), 333--355,
 \href{http://arxiv.org/abs/0709.3844}{arXiv:0709.3844}.

\bibitem{bn.simpleconf}
Bader U., Nevo A., Conformal actions of simple {L}ie groups on compact
 pseudo-{R}iemannian manifolds,
 \href{https://doi.org/10.4310/jdg/1090351122}{\textit{J.~Differential Geom.}}
 \textbf{60} (2002), 355--387.

\bibitem{BFGK}
Baum H., Friedrich T., Grunewald R., Kath I., Twistors and {K}illing spinors on
 {R}iemannian manifolds, \textit{Teubner-Texte Math.}, Vol.~124, B.G.~Teubner
 Verlagsgesellschaft mbH, Stuttgart, 1991.

\bibitem{CD.BGG}
Calderbank D.M.J., Diemer T., Differential invariants and curved
 {B}ernstein--{G}elfand--{G}elfand sequences,
 \href{https://doi.org/10.1515/crll.2001.059}{\textit{J.~Reine Angew. Math.}}
 \textbf{537} (2001), 67--103,
 \href{http://arxiv.org/abs/math.DG/0001158}{arXiv:math.DG/0001158}.

\bibitem{cemn.cproj}
Calderbank D.M.J., Eastwood M.G., Matveev V.S., Neusser K., C-projective
 geometry, \href{https://doi.org/10.1090/memo/1299}{\textit{Mem. Amer. Math.
 Soc.}} \textbf{267} (2020), v+137~pages,
 \href{http://arxiv.org/abs/1512.04516}{arXiv:1512.04516}.

\bibitem{cap.deformations}
{\v{C}}ap A., Infinitesimal automorphisms and deformations of parabolic
 geometries, \href{https://doi.org/10.4171/JEMS/116}{\textit{J.~Eur. Math.
 Soc.}} \textbf{10} (2008), 415--437,
 \href{http://arxiv.org/abs/math.DG/0508535}{arXiv:math.DG/0508535}.

\bibitem{cap.gover.tractor.calc}
{\v{C}}ap A., Gover A.R., Tractor calculi for parabolic geometries,
 \href{https://doi.org/10.1090/S0002-9947-01-02909-9}{\textit{Trans. Amer.
 Math. Soc.}} \textbf{354} (2002), 1511--1548.

\bibitem{cap.me.parabolictrans}
{\v{C}}ap A., Melnick K., Essential {K}illing fields of parabolic geometries,
 \href{https://doi.org/10.1512/iumj.2013.62.5166}{\textit{Indiana Univ.
 Math.~J.}} \textbf{62} (2013), 1917--1953,
 \href{http://arxiv.org/abs/1208.5510}{arXiv:1208.5510}.

\bibitem{cap.schichl.equiv}
{\v{C}}ap A., Schichl H., Parabolic geometries and canonical {C}artan
 connections,
 \href{https://doi.org/10.14492/hokmj/1350912986}{\textit{Hokkaido Math.~J.}}
 \textbf{29} (2000), 453--505.

\bibitem{cap.slovak.book.vol1}
{\v{C}}ap A., Slov\'ak J., Parabolic geometries.~{I}. {B}ackground and general
 theory, \textit{Math. Surveys Monogr.}, Vol.~154,
 \href{https://doi.org/10.1090/surv/154}{American Mathematical Society},
 Providence, RI, 2009.

\bibitem{css.bgg.seqs}
{\v{C}}ap A., Slov\'ak J., Sou\v{c}ek V., Bernstein--{G}elfand--{G}elfand
 sequences, \href{https://doi.org/10.2307/3062111}{\textit{Ann. of Math.}}
 \textbf{154} (2001), 97--113,
 \href{http://arxiv.org/abs/math.DG/0001164}{arXiv:math.DG/0001164}.

\bibitem{curry.gover.tractors}
Curry S.N., Gover A.R., An~introduction to conformal geometry and tractor
 calculus, with a~view to applications in general relativity, in Asymptotic
 Analysis in General Relativity, \textit{London Math. Soc. Lecture Note Ser.},
 Vol.~443, Cambridge University Press, Cambridge, 2018, 86--170,
 \href{http://arxiv.org/abs/1412.7559}{arXiv:1412.7559}.

\bibitem{dani.borel.density}
Dani S.G., On ergodic quasi-invariant measures of group automorphism,
 \href{https://doi.org/10.1007/BF02761685}{\textit{Israel~J. Math.}}
 \textbf{43} (1982), 62--74.

\bibitem{dfmpz.balanced}
Dumitrescu S., Frances C., Melnick K., Pecastaing V., Zeghib A., A~local
 {L}orentzian {F}errand--{O}bata theorem for conformal vector fields,
 \href{http://arxiv.org/abs/2511.03713}{arXiv:2511.03713}.

\bibitem{eastwood.matveev.proj.tractor}
Eastwood M., Matveev V., Metric connections in projective differential
 geometry, in Symmetries and Overdetermined Systems of Partial Differential
 Equations, \textit{IMA Vol. Math. Appl.}, Vol.~144,
 \href{https://doi.org/10.1007/978-0-387-73831-4_16}{Springer}, New York,
 2008, 339--350, \href{http://arxiv.org/abs/0806.3998}{arXiv:0806.3998}.

\bibitem{frances.degenerescence}
Frances C., D\'eg\'enerescence locale des transformations conformes
 pseudo-riemanniennes, \href{https://doi.org/10.5802/aif.2732}{\textit{Ann.
 Inst. Fourier (Grenoble)}} \textbf{62} (2012), 1627--1669,
 \href{http://arxiv.org/abs/1008.2436}{arXiv:1008.2436}.

\bibitem{fz.simpleconf}
Frances C., Zeghib A., Some remarks on conformal pseudo-{R}iemannian actions of
 simple {L}ie groups,
 \href{https://doi.org/10.4310/MRL.2005.v12.n1.a5}{\textit{Math. Res. Lett.}}
 \textbf{12} (2005), 49--56.

\bibitem{GoverSilhan}
Gover A.R., \v{S}ilhan J., The conformal {K}illing equation on
 forms---prolongations and applications,
 \href{https://doi.org/10.1016/j.difgeo.2007.11.014}{\textit{Differential
 Geom. Appl.}} \textbf{26} (2008), 244--266,
 \href{http://arxiv.org/abs/math.DG/0601751}{arXiv:math.DG/0601751}.

\bibitem{gromov.rgs}
Gromov M., Rigid transformations groups, in G\'eom\'etrie Diff\'erentielle
 ({P}aris, 1986), \textit{Travaux en Cours}, Vol.~33, Hermann, Paris, 1988,
 65--139.

\bibitem{hsss.prol.connxn}
Hammerl M., Somberg P., Sou\v{c}ek V., \v{S}ilhan J., On a~new normalization
 for tractor covariant derivatives,
 \href{https://doi.org/10.4171/JEMS/349}{\textit{J.~Eur. Math. Soc.}}
 \textbf{14} (2012), 1859--1883,
 \href{http://arxiv.org/abs/1003.6090}{arXiv:1003.6090}.

\bibitem{kobayashi.transf}
Kobayashi S., Transformation groups in differential geometry,
 \textit{Classics Math.},
 \href{https://doi.org/10.1007/978-3-642-61981-6}{Springer}, Berlin,
 1995.

\bibitem{kostant.harmonic.curv}
Kostant B., Lie algebra cohomology and the generalized {B}orel--{W}eil theorem,
 \href{https://doi.org/10.2307/1970237}{\textit{Ann. of Math.}} \textbf{74}
 (1961), 329--387.

\bibitem{me.frobenius}
Melnick K., A~{F}robenius theorem for {C}artan geometries, with applications,
 \href{https://doi.org/10.4171/LEM/57-1-3}{\textit{Enseign. Math.}}
 \textbf{57} (2011), 57--89,
 \href{http://arxiv.org/abs/0812.0624}{arXiv:0812.0624}.

\bibitem{me.notices}
Melnick K., Rigidity of transformation groups in differential geometry,
 \href{https://doi.org/10.1090/noti2279}{\textit{Notices Amer. Math. Soc.}}
 \textbf{68} (2021), 721--732,
 \href{http://arxiv.org/abs/2009.13937}{arXiv:2009.13937}.

\bibitem{mp.confdambra}
Melnick K., Pecastaing V., The conformal group of a~compact simply connected
 {L}orentzian manifold,
 \href{https://doi.org/10.1090/jams/976}{\textit{J.~Amer. Math. Soc.}}
 \textbf{35} (2022), 81--122,
 \href{http://arxiv.org/abs/1911.06251}{arXiv:1911.06251}.

\bibitem{morimoto}
Morimoto T., Geometric structures on filtered manifolds,
 \href{https://doi.org/10.14492/hokmj/1381413178}{\textit{Hokkaido Math.~J.}}
 \textbf{22} (1993), 263--347.

\bibitem{morris.ratners.thms.book}
Morris D.W., Ratner's theorems on unipotent flows, \textit{Chicago Lectures in Math.},
 University of Chicago Press, Chicago, IL, 2005.

\bibitem{neuss.prolong}
Neusser K., Prolongation on regular infinitesimal flag manifolds,
 \href{https://doi.org/10.1142/S0129167X11007501}{\textit{Internat.~J. Math.}}
 \textbf{23} (2012), 1250007, 41~pages,
 \href{http://arxiv.org/abs/1012.1686}{arXiv:1012.1686}.

\bibitem{pecastaing.frobenius}
Pecastaing V., On two theorems about local automorphisms of geometric
 structures, \href{https://doi.org/10.5802/aif.3009}{\textit{Ann. Inst.
 Fourier (Grenoble)}} \textbf{66} (2016), 175--208,
 \href{http://arxiv.org/abs/1402.5048}{arXiv:1402.5048}.

\bibitem{pecastaing.smooth.sl2r}
Pecastaing V., Lorentzian manifolds with a~conformal action of~{${\rm
 SL}(2,{\bf R})$}, \href{https://doi.org/10.4171/CMH/439}{\textit{Comment.
 Math. Helv.}} \textbf{93} (2018), 401--439.

\bibitem{pecastaing.rank1}
Pecastaing V., Conformal actions of real-rank~1 simple {L}ie groups on
 pseudo-{R}iemannian manifolds,
 \href{https://doi.org/10.1007/s00031-019-09527-6}{\textit{Transform. Groups}}
 \textbf{24} (2019), 1213--1239.

\bibitem{rosenlicht.brasil}
Rosenlicht M., A~remark on quotient spaces, \textit{An. Acad. Brasil. Ci.}
 \textbf{35} (1963), 487--489.

\bibitem{Semmelmann}
Semmelmann U., Conformal {K}illing forms on {R}iemannian manifolds,
 {H}abilitation Thesis, Maximilian Ludwig Universit\"at in M\"unchen, 2001.

\bibitem{shalom.discompact}
Shalom Y., Invariant measures for algebraic actions, {Z}ariski dense subgroups
 and {K}azhdan's property~({T}),
 \href{https://doi.org/10.1090/S0002-9947-99-02363-6}{\textit{Trans. Amer.
 Math. Soc.}} \textbf{351} (1999), 3387--3412.

\bibitem{sharpe}
Sharpe R.W., Differential geometry. {C}artan's generalization of {K}lein's
 {E}rlangen program, \textit{Grad. Texts in Math.}, Vol.~166, Springer, New
 York, 1997.

\bibitem{tanaka}
Tanaka N., On the equivalence problems associated with simple graded {L}ie
 algebras, \href{https://doi.org/10.14492/hokmj/1381758416}{\textit{Hokkaido
 Math.~J.}} \textbf{8} (1979), 23--84.

\bibitem{thurston.reeb.stability}
Thurston W.P., A~generalization of the {R}eeb stability theorem,
 \href{https://doi.org/10.1016/0040-9383(74)90025-1}{\textit{Topology}}
 \textbf{13} (1974), 347--352.

\bibitem{zimmer.etsg}
Zimmer R.J., Ergodic theory and semisimple groups, \textit{Monogr. Math.},
 Vol.~81, \href{https://doi.org/10.1007/978-1-4684-9488-4}{Birkh\"auser},
 Basel, 1984.

\bibitem{zimmer.lorentz}
Zimmer R.J., On the automorphism group of a~compact {L}orentz manifold and
 other geometric manifolds,
 \href{https://doi.org/10.1007/BF01394415}{\textit{Invent. Math.}} \textbf{83}
 (1986), 411--424.

\bibitem{zimmer.rankbounds}
Zimmer R.J., Split rank and semisimple automorphism groups of {$G$}-structures,
 \href{https://doi.org/10.4310/jdg/1214441180}{\textit{J.~Differential Geom.}}
 \textbf{26} (1987), 169--173.

\bibitem{zimmer.gromovrep}
Zimmer R.J., Automorphism groups and fundamental groups of geometric manifolds,
 in Differential Geometry: {R}iemannian Geometry ({L}os {A}ngeles, {CA},
 1990), \textit{Proc. Sympos. Pure Math.}, Vol.~54,
 \href{https://doi.org/10.1090/pspum/054.3/1216656}{American Mathematical
 Society}, Providence, RI, 1993, 693--710.

\end{thebibliography}
\end{document}
